\scanpage{4}
$n$-polyèdres réguliers.

Considérons \deleted{\unsure{le|la}} une $n$-carte régulière, défini par un espace homogène $\underline{R}$ \inserted{(\unsure{tt plein})} sous le groupe $n$-cartographique
\[
\begin{aligned}
G_n=\{\sigma_0,\ldots,\sigma_n\mid {}&\sigma_0^2=\cdots=\sigma_n^2=1\\
&\sigma_i\sigma_j=\sigma_j\sigma_i\quad\text{si }\left\{\begin{aligned}&\deleted{\unsure{i< i+j}}\inserted{j\geq i+2}\\&0\leq i,j\leq n\end{aligned}\right.\}
\end{aligned}
\]
tel que les stabilisateurs \deleted{de}\inserted{$G_{n\underline{x}}$} $\forall\,\underline{x}\in\underline{R}$ sont distingués dans $G_n$, donc indépendants de $x$, soit $H$, \inserted{et $G_n/H=\Gamma$} On suppose que pour tte partie $\unsure{J|I}$ de $[0,n]\cap\mathbb{N}$ ne contenant pas deux entiers consécutifs, le produit \deleted{de} $\sigma_J=\prod_{i\in J}\sigma_i\notin H$ i.e. le produit $\sigma_H$ des $\sigma_i$ \uline{dans} $\Gamma$ est $\neq1$.
% slanted
\margin{inutile \dots}

On pose
\[
\left\{\begin{aligned}
p_1&=\text{ordre de }\sigma_0\sigma_1\\
&\cdots\\
p_n&=\text{ordre de }\sigma_{n-1}\sigma_n
\end{aligned}\right.
\]
\margin{\uline{NB} les $p_i$ peuvent être infinis !}
et on suppose
\[
3\leq p_i\leq+\infty\qquad\text{pour }1\leq i\leq n
\]
\margin{i)}
On suppose ici que $\Gamma$ est le groupe engendré par les $\sigma_i$, avec (en plus des relations plus haut définissant $G_n$) les relations supplémentaires
\[
(\sigma_{i-1}\sigma_i)^{p_i}=1
\]
% archive sheet number: 3
